\begin{ouverture}[calcul]{3}{Puissances et notation scientifique}{Partie 1 · Nombres et calculs}
{Produits répétés · Puissances de dix · Préfixes}
\objectif{Interpréter une puissance à partir de sa définition.}
\objectif{Passer entre écriture décimale et scientifique.}
\objectif{Comparer des ordres de grandeur et utiliser les préfixes.}
\prerequis{Multiplication, division et rangs décimaux.}
\end{ouverture}
\pagestyle{otMain}\essentiel\fontsize{11}{15}\selectfont
\phantomsection\label{cours-3-1}
\essSection{1}{Une écriture de produits}
\begin{definition}Pour un entier $n\ge1$, $a^n$ désigne le produit de $n$ facteurs égaux à $a$. Par exemple $3^4=3\times3\times3\times3=81$. L’exposant indique un nombre de facteurs, pas une multiplication par $n$.\end{definition}
\begin{exemple}$(-2)^4=16$ tandis que $-2^4=-16$ : dans la seconde écriture, la puissance porte seulement sur 2. Pour une base quelconque, ce chapitre utilise des exposants entiers strictement positifs.\end{exemple}
\begin{definition}$10^n$ vaut 1 suivi de $n$ zéros si $n>0$. $10^0=1$ et $10^{-n}=1/10^n$ : $10^{-3}=0{,}001$. Multiplier par $10^n$ ou $10^{-n}$ déplace le rang de chaque chiffre.\end{definition}
\begin{savoirfaire}[Retrouver plutôt que réciter]$10^3\times10^2$ contient cinq facteurs 10, donc vaut $10^5$. $10^4/10^6=1/10^2=10^{-2}$. La définition permet de justifier les calculs sans apprendre une liste de formules générales.\end{savoirfaire}
\clearpage
\phantomsection\label{cours-3-2}
\essSection{2}{Échelles et écriture scientifique}
\begin{definition}La notation scientifique d’un nombre positif non nul est $a\times10^n$, où $1\le a<10$ et $n$ est entier. Par exemple $52\,000=5{,}2\times10^4$ et $0{,}00052=5{,}2\times10^{-4}$. Pour un nombre négatif, on place le signe moins devant l’écriture de sa distance à zéro : $-52\,000=-5{,}2\times10^4$.\end{definition}
\begin{savoirfaire}[Comparer]Comparer d’abord les exposants de deux écritures scientifiques positives, puis les coefficients si les exposants sont égaux. Pour une estimation, arrondir le coefficient à un chiffre significatif et conserver l’unité.\end{savoirfaire}
\begin{center}\begin{tabular}{c|rrrrr}\toprule
Préfixe & nano & micro & milli & centi & déci\\\midrule
Facteur & $10^{-9}$ & $10^{-6}$ & $10^{-3}$ & $10^{-2}$ & $10^{-1}$\\\midrule
Préfixe & déca & hecto & kilo & méga & giga\\\midrule
Facteur & $10^1$ & $10^2$ & $10^3$ & $10^6$ & $10^9$\\\bottomrule
\end{tabular}\end{center}
\begin{exemple}$4{,}5$ mg $=4{,}5\times10^{-3}$ g. Avec le préfixe décimal giga, 2 Go désignent $2\times10^9$ octets ; on n’utilise pas ici les unités binaires.\end{exemple}
\begin{attention}Une écriture $52\times10^3$ est correcte mais n’est pas scientifique, car $52\ge10$. Toujours préciser l’unité quand on compare deux tailles.\end{attention}

\clearpage\rubrique{Découverte}
\begin{blocLivre}{ex-3-1}\begin{exercice}[1]{Déplier une puissance}
Écrire $5^3$ comme produit et le calculer.
\end{exercice}\end{blocLivre}
\begin{blocLivre}{ex-3-2}\begin{exercice}[1]{Attention aux parenthèses}
Calculer $(-3)^2$, $-3^2$ et $(-3)^3$.
\end{exercice}\end{blocLivre}
\begin{blocLivre}{ex-3-3}\begin{exercice}[1]{Puissances de dix}
Donner les écritures décimales de $10^5$, $10^{-4}$ et $10^0$.
\end{exercice}\end{blocLivre}
\begin{blocLivre}{ex-3-4}\begin{exercice}[1]{Vers la notation scientifique}
Écrire $7\,240\,000$ et $0{,}000091$ en notation scientifique.
\end{exercice}\end{blocLivre}
\begin{blocLivre}{ex-3-5}\begin{exercice}[1]{Lire un préfixe}
Convertir 6 km en mètres et 8 microlitres en litres.
\end{exercice}\end{blocLivre}

\rubrique{Entraînement}
\begin{blocLivre}{ex-3-6}\begin{exercice}[2]{Comparer des petites valeurs}
Comparer $4{,}8\times10^{-5}$ et $7{,}2\times10^{-6}$.
\end{exercice}\end{blocLivre}
\begin{blocLivre}{ex-3-7}\begin{exercice}[2]{Retrouver les facteurs dix}
Écrire comme puissance de 10 : $10^4\times10^3$ et $10^2/10^5$.
\end{exercice}\end{blocLivre}
\begin{blocLivre}{ex-3-8}\begin{exercice}[2]{Normaliser une écriture}
Mettre $48\times10^6$ et $0{,}62\times10^{-3}$ en notation scientifique.
\end{exercice}\end{blocLivre}
\begin{blocLivre}{ex-3-9}\begin{exercice}[2]{Multiplier deux écritures}
Calculer $(3\times10^4)(2\times10^{-6})$ et donner une écriture décimale.
\end{exercice}\end{blocLivre}
\begin{blocLivre}{ex-3-10}\begin{exercice}[2]{Diviser et estimer}
Calculer $\dfrac{8{,}4\times10^7}{2\times10^3}$.
\end{exercice}\end{blocLivre}

\rubrique{Problèmes}
\begin{blocLivre}{ex-3-11}\begin{exercice}[3]{Des fichiers}
Une archive contient 2400 images de 2,5 Mo chacune. On utilise les préfixes décimaux. Quelle taille totale en Mo, puis en Go ?
\end{exercice}\end{blocLivre}
\begin{blocLivre}{ex-3-12}\begin{exercice}[3]{Des longueurs microscopiques}
Deux éléments d’un modèle mesurent 450 nm et 0,8 micromètre. Lequel est le plus long et de combien de nanomètres ?
\end{exercice}\end{blocLivre}
\begin{blocLivre}{ex-3-13}\begin{exercice}[3]{Un ordre de grandeur}
Un trajet fictif mesure $3{,}8\times10^5$ km ; on le partage en étapes de $2\times10^3$ km. Combien d’étapes complètes cela représente-t-il ? Estimer d’abord.
\end{exercice}\end{blocLivre}
\begin{blocLivre}{ex-3-14}\begin{exercice}[3]{Une répétition qui double}
Une quantité vaut 3 au départ et double à chaque étape. Exprimer sa valeur après six étapes avec une puissance, puis la calculer.
\end{exercice}\end{blocLivre}
\begin{blocLivre}{ex-3-15}\begin{exercice}[3]{Une somme à ne pas confondre}
Calculer $A=3{,}2\times10^5+8\times10^4$, puis $B=3{,}2\times10^5\times8\times10^4$. Expliquer la différence.
\end{exercice}\end{blocLivre}

\clearpage
